\section{Limitations and Threats to Validity}
\label{sec:limits}

The single largest limitation is stated first and without hedging: \emph{no
external organization has deployed this system}. Every result in this paper was
produced by its authors, on synthetic fixtures or a shared public demonstration
instance, in a controlled environment. There are zero external pilots, zero
production workloads, and therefore zero field-reliability evidence of any kind.
Nothing in this paper should be read as a claim about how the system behaves in
someone else's hands.

The evidence does not establish long-term reliability, maintenance effort,
total cost of ownership, multi-tenant scale, or safety in clinical production.
OpenEMR uses a shared public demonstration site and a small sample. MockMed is a
fixture, and the second-domain lending study is likewise a bounded mechanism
study (twelve tasks, three trials each) on a purpose-built MockLoan ledger, not a
production lending system. The two multi-application benchmarks add branching
and cross-system effects but are equally synthetic, localhost-only, and
hand-authored, with three runs per cell. The breadth corpus has one trial per
application and selection bias toward public no-auth browser apps. No published
study covers weeks of natural drift or a representative distribution of customer
workflows.

Two evaluation paths are measured separately and neither covers the other. The
end-to-end effect study, the lending study, and the two multi-application
benchmarks all actuate through the replayer's API tier with a null observation
backend and a null vision stack, so they measure governed decision, verification,
refusal, and transaction classification---not GUI perception, and not document
understanding. The browser-and-OCR fault study and the public-web breadth corpus
exercise perception, but the fault study carries only a screen-only arm. A study
that exercises real perception \emph{and} effect verification on the same runs
does not yet exist, and no number in this paper should be read as though one
did.

Identity is only checked on armed steps. Risk classification uses keywords and
can miss an icon-only or indirect write. Visual postconditions do not verify a
system-of-record effect. Effect declarations and verifiers require deployment
work. Optional model rescue can false-rescue an ambiguous state, and measured
grounding is weak on dense lists. Human review can catch contextual sanitation
errors but is not proof of de-identification.

Several mechanisms described in Section~\ref{sec:governance} are presented as
mechanisms and are not the subject of a reported measurement. The
transaction-outcome taxonomy is exercised by the multi-application benchmarks,
where it distinguishes \textsc{verified}, \textsc{halted-before-effect}, and
\textsc{reconciliation-required} on hand-authored exception paths, and the
idempotency ledger is exercised by the duplicate-attempt checks in those runs;
but the verifier plugin seam, the governed repair promotion lifecycle, and
explicit surface selection are described from their implementations and their
unit-level contracts, with no benchmark that quantifies, for example, how often
a contract-invariant check correctly refuses a real drifted bundle or how much
operator time promotion costs. Describing a governance mechanism is not
evidence that it holds up under an adversarial operator or at fleet scale.

The substrate qualifications cover one Windows UIA workflow and VM snapshot,
one macOS host and TextEdit task, and one RDP task and snapshot. They establish
working observation, actuation, refusal, and effect-oracle mechanisms in those
environments, not reliability across application families or deployment fleets.
Remote display introduces compression, scaling, latency, coordinate,
lock-screen, credential, input-acceptance, and identity-availability variables.
No measured result covers real Citrix ICA/HDX, so RDP evidence does not transfer
to a Citrix deployment. The Citrix campaign reported in
Section~\ref{sec:results} is a \emph{deterministic synthetic stand-in}: the
shipping backend is real and unmodified, but its window-client seam is replaced
by in-process synthetic frames, so no HDX codec, ICA compression, or
Workspace-client transport is exercised and the effect oracle is an in-process
store. Its degradation conditions are image transforms, not a Thinwire
bitstream. It is evidence about the backend's gating logic under adverse
remote-display conditions and about nothing else; the released artifact and its
status manifest say so in machine-readable form. Citrix is a pixel-only
remote-display environment in the same
class as the measured RDP backend, and the resolution ladder and identity gate
are substrate-independent by construction, but the paper reports no
real-protocol ICA/HDX
measurement. A concrete validation would record and compile a workflow inside a
real Citrix Workspace session, replay it through the vision-only rungs of the
ladder over the ICA/HDX channel, and confirm the identity gate refuses ambiguous
targets and that an out-of-band effect oracle catches an injected transactional
fault. That study is not reported here.

Effect verification also has a substrate-dependent reach that the RDP row can
obscure. The strong oracle is only available when the system of record can be
read out of band, through a database, API, or file interface independent of the
substrate that carried the write. The RDP qualification reads the persisted file
through guest tools, an out-of-band channel that a locked-down Citrix session may
not expose. Where no such read path exists, the transactional strong oracle is
unavailable; a same-application re-read cannot catch a partial, duplicate, stale,
or lost-update write, and on such substrates safety rests on the identity gate
and halt-on-ambiguity rather than on effect verification. Even where a read path
exists, an out-of-band effect oracle is only as strong as its read coverage: in
the end-to-end study a collateral write to a billing surface the oracle did not
read slipped through (9 of 90) until the read path was extended to every mutable
surface, so effect verification guarantees only what it reads.

The same read-coverage limit binds the independent \emph{ground truth}, and this
is why the complete-read-path 0 of 90 must be read as zero in a closed world, not
an absolute zero. The judge audits every persisted table it discovers in the
SQLite system of record, so it is open-world over that database; but an effect
that lands outside the database entirely---an outbound HL7 or message-queue
publish, a filesystem side channel, a downstream service call---is invisible to
both the complete read path and the judge, because no in-database audit can see
it. Independence of code and read path is moreover not independence of
\emph{specification}: the judge and the effect contract encode the same business
intent, so a fault class no one thought to define is invisible to all three
paths. For these reasons the realistic figure we foreground is the 9 of 90
residual under a single out-of-band oracle, the deployment an operator actually
ships; the 0 of 90 is the best case under a complete in-database read path.

Cost values use recorded provider usage and then-current price metadata. They do
not include recording, compilation, operator review, exception handling,
infrastructure, support, or maintenance. Provider prices and models change.

Finally, the authors develop the evaluated system. Released raw artifacts where
the source boundary permits them, bounded aggregate summaries, and scripts
reduce but do not eliminate experimenter bias. Independent replication and
pre-registered longitudinal workloads remain necessary.
